% part: intuitionistic-logic
% chapter: propositions-as-types
% section: terms-to-proofs

\documentclass[../../../include/open-logic-section]{subfiles}

\begin{document}

\olfileid{int}{pty}{tp}

\olsection{सिद्धता-पदांवरून \usetoken{P}{derivation} पुन्हा मिळवणे}

आता उलट दिशा पाहू: सिद्धता-पद~$N$ वरून
\Log{N2i} मधील !!a{proof} पुन्हा मिळवणे.
सर्वसाधारणपणे हे दिलेल्या सिद्धता-पदात मुक्त असलेल्या
प्रत्येक चल~$x$ साठी $x : !A$ अशी जोडी धरणाऱ्या
संदर्भाच्या सापेक्षच करता येते. पण आधी मुक्त चले
नसलेले उदाहरण घेऊ, म्हणजे हे पद:
\[
\lambd[\typeof{x}{!A \land
  !B}][\pair{\proj{2}{x}}{\proj{1}{x}}]
\]
याचे रूप $\lambd[\typeof{x}{!A \land !B}][N_1]$ आहे.
अशा रूपाचे पद निर्माण करणारा \Log{tN2i} चा
एकमेव नियम~\Intro{\lif} असल्यामुळे~$N$ ने
सांकेतिक केलेली !!{proof}~\Intro{\lif} वर संपली
पाहिजे, म्हणजे तिचा शेवट असा आहे:
\[
  \Axiom$x : !A\land !B \fCenter \pair{\proj{2}{x}}{\proj{1}{x}} :
    !C$
  \RightLabel{\Intro\lif}
  \UnaryInf$\fCenter \lambd[\typeof{x}{!A \land
  !B}][\pair{\proj{2}{x}}{\proj{1}{x}}] : (!A \land !B) \lif !C$
  \DisplayProof
  \]
$!C$ काय आहे हे अद्याप माहीत नाही; पण पूर्वांग
$!A \land !B$ आहे, आणि आधारविधानाच्या संदर्भात
$x:!A \land !B$ असलेच पाहिजे हे माहीत आहे.

आता आधारविधानाशी संबंधित सिद्धता-पद निश्चित झाले आहे:
$\pair{\proj{2}{x}}{\proj{1}{x}}$.
\emph{त्याने} सांकेतिक केलेली !!{proof}~\Intro{\land}
वर संपली पाहिजे. ~$!C$ काय आहे हे अजून माहीत
नाही; पण तो~\Intro{\land} चा निष्कर्ष असल्यामुळे
संयोग असलाच पाहिजे, म्हणजे
$!C \ident (!C_1 \land !C_2)$. पुनर्रचना पुढे चालू ठेवू:
\[
  \Axiom$x : !A\land !B \fCenter \proj{2}{x} : !C_1$
  \Axiom$x : !A\land !B \fCenter \proj{1}{x} : !C_2$
  \RightLabel{\Intro\land}
  \BinaryInf$x : !A\land !B \fCenter \pair{\proj{2}{x}}{\proj{1}{x}} :
    !C_1 \land !C_2$
  \RightLabel{\Intro\lif}
  \UnaryInf$\fCenter \lambd[\typeof{x}{!A \land
  !B}][\pair{\proj{2}{x}}{\proj{1}{x}}] : (!A \land !B) \lif (!C_1 \land !C_2)$
  \DisplayProof
  \]
$\pi_2(x)$ हे सिद्धता-पद सांगते की डावीकडील वरची
क्रमवर्ती~\Elim{\land}[2] ने मिळाली; म्हणजे तिचे
आधारविधान $!D \land !C_1$ या रूपाचे आहे. उजवीकडे
$!C_2$ कडे नेणारे अनुमान~\Elim{\land}[1] असले
पाहिजे आणि त्याचे आधारविधान $!C_2 \land !E$ या
रूपाचे असले पाहिजे, हे ठरवता येते.
\[
  \Axiom$x : !A\land !B \fCenter x: !D\land !C_1$
  \RightLabel{\Elim\land[2]}
  \UnaryInf$x : !A\land !B \fCenter \proj{2}{x} : !C_1$
  \Axiom$x : !A\land !B \fCenter x: !C_2\land !E$
  \RightLabel{\Elim\land[1]}
  \UnaryInf$x : !A\land !B \fCenter \proj{1}{x} : !C_2$
  \RightLabel{\Intro\land}
  \BinaryInf$x : !A\land !B \fCenter \pair{\proj{2}{x}}{\proj{1}{x}} :
    !C_1 \land !C_2$
  \RightLabel{\Intro\lif}
  \UnaryInf$\fCenter \lambd[\typeof{x}{!A \land
  !B}][\pair{\proj{2}{x}}{\proj{1}{x}}] : (!A \land !B) \lif (!C_1 \land !C_2)$
  \DisplayProof
  \]
सर्वांत वरच्या क्रमवर्तींची सिद्धता-पदे आता चले
असल्यामुळे त्या स्वयंसिद्धकेच असल्या पाहिजेत.
त्यावरून $!C_1$, $!C_2$, $!D$, आणि $!E$ ठरतात:
$!D \ident !A$ आणि $!C_1 \ident !B$ नसतील,
तर डावी क्रमवर्ती स्वयंसिद्धक नाही; तसेच
$!C_2 \ident !A$ आणि $!E \ident !B$ नसतील,
तर उजवी क्रमवर्तीही स्वयंसिद्धक नाही.
अशा रीतीने संपूर्ण सिद्धता पुन्हा मिळाली:
\[
  \Axiom$x : !A\land !B \fCenter x: !A\land !B$
  \RightLabel{\Elim\land[2]}
  \UnaryInf$x : !A\land !B \fCenter \proj{2}{x} : !B$
  \Axiom$x : !A\land !B \fCenter x: !A\land !B$
  \RightLabel{\Elim\land[1]}
  \UnaryInf$x : !A\land !B \fCenter \proj{1}{x} : !A$
  \RightLabel{\Intro\land}
  \BinaryInf$x : !A\land !B \fCenter \pair{\proj{2}{x}}{\proj{1}{x}} :
    !B \land !A$
  \RightLabel{\Intro\lif}
  \UnaryInf$\fCenter \lambd[\typeof{x}{!A \land
  !B}][\pair{\proj{2}{x}}{\proj{1}{x}}] : (!A \land !B) \lif (!B \land !A)$
  \DisplayProof
  \]

\end{document}
